\documentclass[../../main.tex]{subfiles}% 注意这里的文件路径不能用 ./main.tex ,否则用latexmk编译子文件会报错
\graphicspath{{\subfix{./image/}}} % 指定图片目录,后续可以直接使用图片文件名
% 注意这里的文件路径不能用 ../../image/ ,否则用latexmk编译子文件会报错

% 例如:
% \begin{figure}[H]
% \centering
% \includegraphics[scale=0.3]{图.png}
% \caption{}
% \label{figure:图}
% \end{figure}
% 注意:上述\label{}一定要放在\caption{}之后,否则引用图片序号会只会显示??.

\begin{document}

\section{正则参数曲面}
在三维欧氏空间$\mathbb{E}^3$中，正则参数曲面是由区域$D\subset\mathbb{R}^2$到$\mathbb{E}^3$的映射$\bm{r}(u,v)$构成的，满足$\bm{r}(u,v)$是三次以上连续可微函数，且$\bm{r}_u\times\bm{r}_v\neq\bm{0}$，其中$\bm{r}_u=\dfrac{\partial\bm{r}}{\partial u}$，$\bm{r}_v=\dfrac{\partial\bm{r}}{\partial v}$，$(u,v)$是曲面的曲纹坐标。

\begin{definition}
若映射$\bm{r}:D\to\mathbb{E}^3$满足
\begin{enumerate}[(1)]
\item $\bm{r}(u,v)=(x(u,v),y(u,v),z(u,v))$是三次以上连续可微函数；
\item $\bm{r}_u\times\bm{r}_v\neq\bm{0},\ \forall(u,v)\in D$，
\end{enumerate}
则称$\bm{r}(u,v)$为$\mathbb{E}^3$中的**正则参数曲面**，$\bm{r}_u,\bm{r}_v$线性无关的点称为曲面的正则点。
\end{definition}

容许参数变换满足
\[
\begin{cases}
u=u(\tilde{u},\tilde{v}),\\
v=v(\tilde{u},\tilde{v})
\end{cases}
\]
是三次以上连续可微函数，且Jacobi行列式$\dfrac{\partial(u,v)}{\partial(\tilde{u},\tilde{v})}\neq0$，变换后仍为正则参数曲面，且
\[
\bm{r}_{\tilde{u}}\times\bm{r}_{\tilde{v}}=\dfrac{\partial(u,v)}{\partial(\tilde{u},\tilde{v})}\cdot(\bm{r}_u\times\bm{r}_v).
\]

\begin{example}
圆柱面$\bm{r}(u,v)=(a\cos u,a\sin u,bv)$（$a>0$）、旋转面$\bm{r}(u,v)=(f(v)\cos u,f(v)\sin u,g(v))$、正螺旋面$\bm{r}(u,v)=(u\cos v,u\sin v,av)$均为正则参数曲面。
\end{example}

\begin{example}
直纹面的参数方程为$\bm{r}(u,v)=\bm{a}(u)+v\bm{l}(u)$，其中$\bm{a}(u)$是准线，$\bm{l}(u)$是直母线方向向量，正则的充要条件是$(\bm{a}'(u)+v\bm{l}'(u))\times\bm{l}(u)\neq\bm{0}$。
\end{example}

\section{切平面和法线}
\begin{definition}
曲面$\bm{r}(u,v)$上经过点$p$的任意连续可微曲线在该点的切向量，称为曲面在点$p$的切向量，全体切向量构成的二维线性空间称为**切空间**，记作$T_pS$。
\end{definition}

切空间由$\bm{r}_u,\bm{r}_v$张成，切平面方程为
\[
\bm{X}=\bm{r}(u,v)+\lambda\bm{r}_u+\mu\bm{r}_v,\quad \lambda,\mu\in\mathbb{R}.
\]
单位法向量为
\[
\bm{n}(u,v)=\dfrac{\bm{r}_u\times\bm{r}_v}{|\bm{r}_u\times\bm{r}_v|},
\]
法线方程为
\[
\bm{X}(t)=\bm{r}(u,v)+t\bm{n}(u,v),\quad t\in\mathbb{R}.
\]

若曲面是等值面$f(x,y,z)=c$，则$\left(\dfrac{\partial f}{\partial x},\dfrac{\partial f}{\partial y},\dfrac{\partial f}{\partial z}\right)$是曲面的法向量。

\section{第一基本形式}
\begin{definition}
曲面的第一类基本量为
\[
E=\bm{r}_u\cdot\bm{r}_u,\quad F=\bm{r}_u\cdot\bm{r}_v,\quad G=\bm{r}_v\cdot\bm{r}_v,
\]
二次微分形式
\[
\mathrm{I}=E(\mathrm{d}u)^2+2F\mathrm{d}u\mathrm{d}v+G(\mathrm{d}v)^2
\]
称为曲面的**第一基本形式**。
\end{definition}

第一基本形式是切向量长度的度量，即$|\mathrm{d}\bm{r}|^2=\mathrm{I}$，其中$\mathrm{d}\bm{r}=\bm{r}_u\mathrm{d}u+\bm{r}_v\mathrm{d}v$。

两个切向量$\mathrm{d}\bm{r},\delta\bm{r}$的内积为
\[
\mathrm{d}\bm{r}\cdot\delta\bm{r}=E\mathrm{d}u\delta u+F(\mathrm{d}u\delta v+\mathrm{d}v\delta u)+G\mathrm{d}v\delta v,
\]
夹角余弦为
\[
\cos\angle(\mathrm{d}\bm{r},\delta\bm{r})=\dfrac{E\mathrm{d}u\delta u+F(\mathrm{d}u\delta v+\mathrm{d}v\delta u)+G\mathrm{d}v\delta v}{\sqrt{E(\mathrm{d}u)^2+2F\mathrm{d}u\mathrm{d}v+G(\mathrm{d}v)^2}\sqrt{E(\delta u)^2+2F\delta u\delta v+G(\delta v)^2}}.
\]

\begin{theorem}
曲面上参数曲线网为正交曲线网的充要条件是$F\equiv0$。
\end{theorem}

\begin{proof}
参数曲线切向量为$\bm{r}_u,\bm{r}_v$，夹角余弦为$\dfrac{F}{\sqrt{EG}}$，故正交等价于$F=0$。

\end{proof}

曲面的面积元素为$\mathrm{d}\sigma=\sqrt{EG-F^2}\mathrm{d}u\mathrm{d}v$，曲面面积为
\[
A=\iint\limits_D \sqrt{EG-F^2}\mathrm{d}u\mathrm{d}v.
\]

\begin{example}
旋转面$\bm{r}(u,v)=(f(v)\cos u,f(v)\sin u,g(v))$的第一基本形式为
\[
\mathrm{I}=f^2(v)(\mathrm{d}u)^2+\left((f'(v))^2+(g'(v))^2\right)(\mathrm{d}v)^2.
\]
\end{example}

\section{曲面上正交参数曲线网的存在性}
\begin{lemma}
设$f(u,v),g(u,v)$是区域$D$上不同时为零的连续可微函数，则对任意$(u_0,v_0)\in D$，存在邻域$U\subset D$及非零连续函数$\lambda(u,v)$，使得$\lambda(f\mathrm{d}u+g\mathrm{d}v)$是全微分。
\end{lemma}

\begin{theorem}
正则参数曲面的每一点邻域内存在正交参数系，使得$F\equiv0$。
\end{theorem}

\begin{proof}
对自然基底$\{\bm{r}_u,\bm{r}_v\}$作Schmidt正交化，得到单位正交标架场$\{\bm{e}_1,\bm{e}_2\}$，由引理可知存在参数系，使得参数曲线切向量与$\bm{e}_1,\bm{e}_2$平行，故为正交参数系。

\end{proof}

\section{保长对应和保角对应}
\begin{definition}
若曲面$S_1$到$S_2$的连续可微映射$\sigma$保持切向量长度不变，即$\sigma^*\mathrm{I}_2=\mathrm{I}_1$，则称$\sigma$为**保长对应**。
\end{definition}

\begin{definition}
若映射$\sigma$保持切向量夹角不变，即$\sigma^*\mathrm{I}_2=\lambda^2\mathrm{I}_1$（$\lambda>0$为比例因子），则称$\sigma$为**保角对应**。
\end{definition}

\begin{theorem}
任意正则参数曲面的每一点邻域，都可与平面区域建立保角对应。
\end{theorem}

\section{可展曲面}
\begin{definition}
直纹面$\bm{r}(u,v)=\bm{a}(u)+v\bm{l}(u)$的高斯曲率恒为零，则称该曲面为**可展曲面**，可展曲面局部是平面、柱面、锥面或切线面。
\end{definition}

\begin{remark}
可展曲面可以无撕裂地展开为平面，即与平面存在保长对应。
\end{remark}




\end{document}